\section{Conclusion}

This paper studies whether LLMs can remove or reduce the manual labeling step in supervised entity matching. The proposed workflow selects candidate pairs from two source tables, labels them with an LLM, and uses the resulting selected-and-labeled data to train a downstream matcher. The preliminary results show that targeted active-learning-based construction can match or outperform official training splits on several benchmarks, while simpler similarity-based selection is less reliable.

The main lesson is that automated EM training-data construction is a selection-and-labeling problem. A capable LLM labeler is necessary, but not sufficient; the system also needs to spend the labeling budget on positives and boundary-relevant negatives. The training-set audit explains this effect by showing that the strongest selected-and-labeled sets differ systematically from simpler retrieval-based sets and from official splits. The remaining experiments should complete the cost table, add a compact robustness check with additional LLM labelers including an open-weight model, and connect the training-set composition analysis more directly to the benchmark-level F1 results.
